% TOP_PROBLEM: {"id": 30007224, "title": "Polynomial-bond PEPS approximation of every gapped square-grid ground vector", "category_id": 16, "category": {"id": 16, "name": "physics", "display_name": "Mathematical Physics", "description": "Problems at the intersection of mathematics and physics.", "slug": "physics", "order_index": 16, "created_at": "2026-07-31T15:26:25.671Z"}, "status": "open"}
% FIRST_POSTED: 2026-10-06
% ATTEMPT_STATUS: solved
% OPENAI_CLAIM: {"schema_version":1,"source_repo":"https://github.com/openai/math","source_commit":"adc7f1241b42e322a6451854ab7e4b4c146bf78a","release_date":"2026-10-06","match":"full","resolution":"proved","independently_reviewed":false,"second_pass":"agreed","families":[{"family":"265","title":"Area laws and tensor networks for two-dimensional gapped systems","subject":"Mathematical physics","manuscripts":[{"title":"Polynomial PEPS approximation of gapped square-grid ground states","pdf_path":"preprints/Polynomial-PEPS-approximation-of-gapped-square-grid-ground-states-September-24-2026/paper.pdf","readme_path":"preprints/Polynomial-PEPS-approximation-of-gapped-square-grid-ground-states-September-24-2026/README.md","date":"2026-09-24","theorem_ref":"thm:main (sections/00-introduction.tex:36)","statement_sources":[{"path":"preprints/Polynomial-PEPS-approximation-of-gapped-square-grid-ground-states-September-24-2026/build/sections/00-introduction.tex","line":36,"label":"thm:main"}]}],"lean":null,"reasoning_trace":null}]}
\documentclass[11pt]{article}
\usepackage[margin=25mm]{geometry}
\usepackage{fontspec}
\setmainfont{Latin Modern Roman}
\usepackage{amsmath,amssymb,mathtools}
\usepackage{unicode-math}
\setmathfont{Latin Modern Math}
\usepackage{enumitem}
\usepackage{xurl,hyperref,bookmark}
\hypersetup{colorlinks=true,urlcolor=blue}
\setcounter{secnumdepth}{0}
\setlength{\parindent}{0pt}
\setlength{\parskip}{5pt}
\setlength{\emergencystretch}{3em}
\begin{document}
\section*{Polynomial-bond PEPS approximation of every gapped square-grid ground vector}

\subsection{OpenAI's claim}
OpenAI presents a proof in the selected manuscript. Uniformly gapped square-grid ground states are claimed to admit PEPS approximations with polynomial bond dimension and normalized vector error at most $L^{-1}$.

\subsection{Mathematical statement of the claimed result}
\paragraph{Result 1.} Fix local dimension \(q\ge2\) and \(J,\Delta>0\). On an \(L\)-by-\(L\) square grid, let
\[
 H=\sum_vh_v+\sum_eh_e,\quad\|h_v\|,\|h_e\|\le J,\qquad
 H-E_0I\ge\Delta(I-|\Omega\rangle\langle\Omega|),
\]
with site/edge-supported Hermitian terms and a unique normalized ground vector \(\Omega\). The claimed constants \(C,c>0\), depending only on \(q,J,\Delta\), give a nonzero projected entangled-pair state \(\Phi\) for every \(L\ge2\), with maximum bond dimension \(CL^c\), satisfying
\[
 \min_{\theta\in\mathbb R}\|\Phi/\|\Phi\|-e^{i\theta}\Omega\|\le L^{-1}.
\]
The assertion is existence of the PEPS, without an efficient construction guarantee.

\subsection{Outline of the argument}
Separated-region information bounds lead to adaptive finite-rank constraints on small patches. Distributed density networks are compressed using encoded frames and controlled ownership changes; a planar distribution protocol assembles the PEPS without assuming gaps for subregions.

\subsection{Scope relative to this problem}
Uniformly gapped square-grid ground states are claimed to admit PEPS approximations with polynomial bond dimension and normalized vector error at most $L^{-1}$. The question retains these same parameters and conventions. Within these assumptions, this is a claimed affirmative answer to the stated question.

\subsection{Formal verification}
Family 265 has no Lean scope page in this release. Its selected result is not presented as formalized. This site has not executed the Lean code. The scope descriptions do not certify the remaining manuscript arguments.

\subsection{Status caveat}
OpenAI attributes these manuscripts to work by an internal model. Its release cautions that arguments without formalization can contain mistakes. This catalogue has not independently checked the proofs; the labels report the mathematical claims attributed to OpenAI.

\subsection{Sources}
\begin{itemize}
\item OpenAI, \emph{Polynomial PEPS approximation of gapped square-grid ground states}, 2026-09-24; citation key \texttt{OAI:Polynomial-PEPS-approximation-of-gapped-square-grid-ground-states-September-24-2026}. \href{https://github.com/openai/math/blob/adc7f1241b42e322a6451854ab7e4b4c146bf78a/preprints/Polynomial-PEPS-approximation-of-gapped-square-grid-ground-states-September-24-2026/paper.pdf}{PDF (GitHub)}; \href{https://github.com/openai/math/raw/adc7f1241b42e322a6451854ab7e4b4c146bf78a/preprints/Polynomial-PEPS-approximation-of-gapped-square-grid-ground-states-September-24-2026/paper.pdf}{PDF (direct)}; \href{https://github.com/openai/math/blob/adc7f1241b42e322a6451854ab7e4b4c146bf78a/preprints/Polynomial-PEPS-approximation-of-gapped-square-grid-ground-states-September-24-2026/README.md}{Manuscript page and citation}; \href{https://github.com/openai/math/blob/main/preprints/Polynomial-PEPS-approximation-of-gapped-square-grid-ground-states-September-24-2026/paper.pdf}{Latest version}; Statement source: \href{https://github.com/openai/math/blob/adc7f1241b42e322a6451854ab7e4b4c146bf78a/preprints/Polynomial-PEPS-approximation-of-gapped-square-grid-ground-states-September-24-2026/build/sections/00-introduction.tex\#L36}{thm:main}.
\item \href{https://github.com/openai/math/blob/adc7f1241b42e322a6451854ab7e4b4c146bf78a/CONTENTS.md}{OpenAI catalogue entry}.
\end{itemize}
\end{document}
